\documentclass[10pt,a4paper]{amsart}
\usepackage[margin=19mm]{geometry}
\usepackage{amsmath,amssymb,mathtools}
\usepackage[scr=rsfs,cal=euler]{mathalfa}
\usepackage{xfrac}
\usepackage[unicode,hidelinks,bookmarks=false]{hyperref}
\newcommand{\ie}{i.e.\ }
\newcommand{\cf}{cf.\ }
\newcommand{\resp}{respectively\ }
\newcommand{\Subsection}{\subsection}
\providecommand{\nobreakdash}{\nobreak\hbox{-}}
\setlength{\emergencystretch}{2em}
\multlinegap0pt
\usepackage{bm}
\usepackage{bbm}
\usepackage{dsfont}
\usepackage{xspace}
\usepackage{cancel}
\usepackage{mathtools}
\usepackage{amsthm}
\makeatletter
\@ifundefined{theorem}{\newtheorem{theorem}{Theorem}}{}
\@ifundefined{lemma}{\newtheorem{lemma}[theorem]{Lemma}}{}
\makeatother
\DeclareMathOperator{\fixedpoint}{\mathfrak K}
\newcommand{\asympint}[2]{\overset {(#1,#2)}\asymp }
\title[Self-similar solutions to the time-fractional Porous-Medium ]{Self-similar solutions to the time-fractional Porous-Medium Equation}
\author{David G\'omez-Castro, {\L}ukasz P{\l}ociniczak, Juan Luis V\'azquez}
\date{Source: arXiv:2604.09281v1 (CC BY 4.0)}
\begin{document}
\maketitle

\noindent\emph{Source and licence.} Self-similar solutions to the time-fractional Porous-Medium Equation. Version arXiv:2604.09281v1 (CC BY 4.0), distributed under the Creative Commons Attribution 4.0 International licence, as recorded in the arXiv licence metadata for this exact version and shown on the abstract page https://arxiv.org/abs/2604.09281v1. Full rights notice: licenses/arxiv-2604.09281-CC-BY-4.0.txt.\par\smallskip

\noindent\textit{Editorial scope.} Counted unit: the Lemma whose source label is \texttt{None} in the article (source0.tex lines 1370--1372, source label \texttt{None}), delivered together with the article's own complete original proof of it (source0.tex lines 1374--1507, one \texttt{proof} environment of 134 lines). No mathematics is written, repaired or re-derived: statement and proof are the article's own text line for line. Required context retained uncounted: eq:Q away from 0 (lines 1101--1104). Every definition, display and label of those lines is retained unchanged. Omitted from this package and not redistributed: 1--1100, 1105--1369, 1373--1373, 1508--3327. Nothing inside a retained excerpt is omitted or altered: every retained body is delivered exactly as the article's lines are written, so no omission receipt of any kind accompanies this package. Citations: the retained excerpts carry no citation command at all, so no bibliography entry of the article is transcribed and no reference is redistributed. Reference adaptation: none was needed and none was made; every reference inside the delivered unit resolves inside it, so no unresolved reference or citation is produced. Printed slips kept literal: no printed word, symbol, hypothesis or number of the counted statement or of its proof was altered. Layout and class: the article's own class is \texttt{article} with packages including fullpage, comment, xcolor, amsthm, amsmath, amssymb, mathtools, hyperref, cleveref, enumitem, todonotes, biblatex, authblk; the wrapper is the lane's pinned \texttt{amsart} editorial wrapper at 10pt on A4, which restates the theorem-like environments the retained text needs and adds the article's own macro definitions that the retained text needs (\texttt{\textbackslash asympint}, \texttt{\textbackslash fixedpoint}), with extra wrapper packages (bm, bbm, dsfont, xspace, cancel, mathtools). The delivered numbering is the unit's own. Assets: no retained span contains a figure environment, a graphics-inclusion command, a PDF-page inclusion, a table or a TikZ picture; no image file is copied into this package, the package declares no asset dependency and no image file is redistributed with it. Licence: version arXiv:2604.09281v1 of this article is distributed under the Creative Commons Attribution 4.0 International licence, as recorded in the arXiv licence metadata for this exact version and shown on the abstract page https://arxiv.org/abs/2604.09281v1. A full rights notice is delivered at licenses/arxiv-2604.09281-CC-BY-4.0.txt. Characters: every retained excerpt is pure ASCII, so the delivered engine, XeLaTeX with its default Latin Modern text font, reports no missing character.
\begin{equation}
  \label{eq:Q away from 0}
  Q(\eta) \asympint{\frac{1}{2}}{1} (1 - \eta)^{1-\alpha}
\end{equation}

\begin{lemma}
Let $\alpha \in (0,1)$, $m > 1$, and $d \ge 1$. We have $\fixedpoint \mathcal V_{\alpha,m} \asymp \mathcal V_{\alpha,m}$.
\end{lemma}

\begin{proof}
We separate the proof into different ranges of $z$. We will show only the proof of the upper bounds, and the lower bounds are proved similarly.

\subparagraph{Range $z \in (\frac 1 4,1)$ and $d \ge 1$.}
We have $K(z,\rho) = \rho Q(z/\rho)$ and hence using \eqref{eq:Q away from 0} we have that using the change of variable $(1-z) \sigma = 1-\rho$
\begin{align*}
  \int_z^1 K(z,\rho) \mathcal V_{\alpha,m} (\rho) d \rho
  & =   \int_z^1 \rho Q(z/\rho) (1-\rho)_+^{\frac{2-\alpha}{m-1}} d \rho
  \\
  & \le  C(\alpha,m,d)   \int_z^1 \rho (1-z/\rho)^{1-\alpha} (1-\rho)_+^{\frac{2-\alpha}{m-1}} d \rho
  \\
  %
  %
  %
  %
  & \le  C(\alpha,m,d) (1-z)^{1-\alpha + \frac{2-\alpha}{m-1} + 1} \int_0^1 (1 - \sigma)^{1-\alpha} \sigma^{\frac{2-\alpha}{m-1}} d \sigma
  %
  %
  \\
  & \le {C(\alpha,m,d) } \mathcal V_{\alpha,m}(z)^m.
\end{align*}
And the lower bound follows similarly bounding $\rho \ge z \ge \frac 1 4$.

\subparagraph{Range $z \in (0,1/4)$ when $d = 1$.}
We have that
\begin{align*}
  \int_z^1 K(z,\rho) \mathcal V (\rho) d \rho
  %
  %
  & \le C \left(  \int_z^{2z}  Q(z/ \rho) \rho d \rho + \int_{2z}^{1/2}   Q(z/\rho) \rho d \rho + \int_{1/2}^{1}  Q(z/\rho) (1-\rho)^{\frac{2-\alpha}{m-1}} d \rho
  \right)
  \\
  & \le C\left( \int_z^{2z}  (1 - z/\rho)^{1-\alpha}  \rho d \rho + \int_{2z}^{1/2}\rho d \rho+ \int_{1/2}^{1} (1-\rho)^{\frac{2-\alpha}{m-1}} d \rho
  \right)
  \\
  & \le C \left(
    z^2
  +  1 \right)
  %
  %
  \le C \mathcal V(z)^m.
\end{align*}
The lower bound follows equivalently.

\subparagraph{Range $z \in (0,1/4)$ when $d \ge 3$.}
We split the computation into three parts
\begin{align*}
  \int_z^1 K(z,\rho) \mathcal V_{\alpha,m} (\rho) d \rho
  & \le C\left( \int_z^{2z}  K(z, \rho) \rho^{-\gamma} d \rho
    + \int_{2z}^{1/2}   K(z, \rho) \rho^{-\gamma} d \rho
  + \int_{1/2}^{1}  K(z, \rho) (1-\rho)^{\frac{2-\alpha}{m-1}} d \rho \right).
\end{align*}
Now we deal with each term separately
\begin{align*}
  \int_z^{2z}  K(z, \rho) \rho^{-\gamma} d \rho
  & \le C \int_z^{2z}  Q(z/ \rho) \rho^{1-\gamma} d \rho
  \le C\int_z^{2z}  (1 - z/\rho)^{1-\alpha}  \rho^{1-\gamma} d \rho
  %
  %
  %
  %
  \le C z^{2-\gamma}.
\end{align*}
We also have that
\begin{align*}
  \int_{2z}^{1/2}   K(z, \rho) \rho^{-\gamma} d \rho
  & \le C \int_{2z}^{1/2}   Q(z/\rho) \rho^{1-\gamma} d \rho
  \le C \int_{2z}^{1/2}(z/\rho)^{2-d}\rho^{1-\gamma} d \rho
  %
  %
  %
  \le C z^{2-d} .
\end{align*}
Lastly
\begin{align*}
  \int_{1/2}^{1}  K(z, \rho) (1-\rho)^{\frac{2-\alpha}{m-1}} d \rho
  & \le C \int_{1/2}^{1}  Q(z/\rho) (1-\rho)^{\frac{2-\alpha}{m-1}} d \rho
  %
  %
  \le C \int_{1/2}^{1} (z/\rho)^{2-d} (1-\rho)^{\frac{2-\alpha}{m-1}} d \rho
  %
  %
  %
  %
  \le Cz^{2-d}.
\end{align*}
Because $m > 1$, we have $\gamma < d-2$, which implies that the term $z^{2-d}$ strictly dominates the term $z^{2-\gamma}$ as $z \to 0^+$. Given that $\gamma m = d-2$, this yields the correct target scaling. We conclude that $(\fixedpoint\mathcal{V}_d(z))^m \le C z^{2-d} \le C \mathcal{V}_d(z)^m$. The lower bound follows analogously.

\subparagraph{Range $z \in (0,1/4)$ when $d = 2$.}
We proceed in a similar way, computing each term
\begin{align*}
  \int_z^1 & K(z,\rho) \mathcal V_{\alpha,m} (\rho) d \rho
  \\
  & \le C \left( \int_z^{2z}  K(z, \rho) (-\log \rho)^{\frac 1 m} d \rho + \int_{2z}^{1/2}   K(z, \rho) (-\log \rho) ^{\frac 1 m}d \rho + \int_{1/2}^{1}  K(z, \rho) (1-\rho)^{\frac{2-\alpha}{m-1}} d \rho \right).
\end{align*}
We now estimate the terms separately. First, we deal with the influence of $\rho$ closest to $z$
\begin{align*}
  \int_z^{2z}  K(z, \rho) (-\log \rho)^{\frac 1 m} d \rho
  & \le C \int_z^{2z}  Q(z/ \rho) \rho(-\log \rho)^{\frac 1 m} d \rho
  \\
  & \le C  \int_z^{2z}  (1 - z/\rho)^{1-\alpha}  \rho(-\log \rho)^{\frac 1 m} d \rho
  \\
  & \le C z^2 \int_1^{2}  (1 - \sigma^{-1})^{1-\alpha}  \sigma (-\log z - \log \sigma)^{\frac 1 m} d \sigma
  \\
  & \le  C z^2 (-\log z)^{\frac 1 m} \int_1^{2}  (1 - \sigma^{-1})^{1-\alpha}  \sigma  d \sigma.
\end{align*}
Now we deal with the influence of intermediate values of $\rho$
\begin{align*}
  \int_{2z}^{1/2}  K(z, \rho) (-\log \rho)^{\frac 1 m} d \rho
  & \le \int_{2z}^{1/2}   Q(z/\rho) \rho(-\log \rho)^{\frac 1 m} d \rho
  %
  %
  \le C \int_{2z}^{1/2}(-\log z/\rho)\rho(-\log \rho)^{\frac 1 m} d \rho
  \\
  & \le C\left( (-\log z) \int_{z}^{\frac 1 2} \rho (-\log \rho)^{\frac 1 m} d\rho -   \int_{z}^{\frac 1 2} \rho (-\log \rho)^{\frac 1 m + 1}d\rho \right)
  \\
  & \le C (-\log z).
\end{align*}
Lastly, we control the terms where $\rho$ is very far from $z$
\begin{align*}
  \int_{\frac 1 2}^{1}  K(z, \rho) (-\log \rho)^{\frac 1 m} d \rho
  & \le \int_{1/2}^{1}  Q(z/\rho) (1-\rho)^{\frac{2-\alpha}{m-1}} d \rho
  %
  %
  \le C \int_{1/2}^{1} (-\log z/\rho) (1-\rho)^{\frac{2-\alpha}{m-1}} d \rho
  \\
  & 
  \le  C (-\log (z/2)) \int_{1/2}^{1} (1-\rho)^{\frac{2-\alpha}{m-1}} d \rho
  %
  %
  \le C (-\log z).
\end{align*}
Since $m > 1$, we conclude that $(\fixedpoint\mathcal V (z))^{m} \le C (-\log z) \le C \mathcal V(z)^m$. 
\end{proof}

\end{document}
